\apendice{Especificación de Requisitos}

\section{Introducción}

En este anexo se recogen los requisitos funcionales y no funcionales definidos para el desarrollo del sistema SIEM. Asimismo, se describen los principales casos de uso que permiten comprender las funcionalidades implementadas y la interacción del usuario con el sistema.

La especificación de requisitos ha servido como guía durante todo el proceso de desarrollo, permitiendo verificar posteriormente que los objetivos planteados inicialmente han sido alcanzados.


Una muestra de cómo podría ser una tabla de casos de uso:

% Caso de Uso 1 -> Consultar Experimentos.
\begin{table}[p]
	\centering
	\begin{tabularx}{\linewidth}{ p{0.21\columnwidth} p{0.71\columnwidth} }
		\toprule
		\textbf{CU-1}    & \textbf{Ejemplo de caso de uso}\\
		\toprule
		\textbf{Versión}              & 1.0    \\
		\textbf{Autor}                & Alumno \\
		\textbf{Requisitos asociados} & RF-xx, RF-xx \\
		\textbf{Descripción}          & La descripción del CU \\
		\textbf{Precondición}         & Precondiciones (podría haber más de una) \\
		\textbf{Acciones}             &
		\begin{enumerate}
			\def\labelenumi{\arabic{enumi}.}
			\tightlist
			\item Pasos del CU
			\item Pasos del CU (añadir tantos como sean necesarios)
		\end{enumerate}\\
		\textbf{Postcondición}        & Postcondiciones (podría haber más de una) \\
		\textbf{Excepciones}          & Excepciones \\
		\textbf{Importancia}          & Alta o Media o Baja... \\
		\bottomrule
	\end{tabularx}
	\caption{CU-1 Nombre del caso de uso.}
\end{table}

\section{Objetivos generales}
Los objetivos generales definidos para el sistema son los siguientes:

\begin{itemize}
\item Centralizar la información procedente de diferentes dispositivos de seguridad.
\item Almacenar de forma persistente los eventos generados por el entorno monitorizado.
\item Detectar situaciones potencialmente peligrosas mediante reglas de correlación.
\item Identificar anomalías relacionadas con accesos, sensores y condiciones ambientales.
\item Facilitar la consulta de información mediante una API REST.
\item Proporcionar una interfaz gráfica que permita supervisar el estado del sistema.
\item Mantener una arquitectura modular que facilite futuras ampliaciones.
\end{itemize}

\section{Catálogo de requisitos}
\subsection{Requisitos funcionales}

\begin{itemize}
\item \textbf{RF-01}: El sistema deberá procesar eventos procedentes de diferentes dispositivos simulados.
\item \textbf{RF-02}: El sistema deberá almacenar los eventos en una base de datos.
\item \textbf{RF-03}: El sistema deberá generar alertas mediante reglas de correlación.
\item \textbf{RF-04}: El sistema deberá detectar anomalías sobre los eventos almacenados.
\item \textbf{RF-05}: El sistema deberá proporcionar una API REST para la consulta de información.
\item \textbf{RF-06}: El sistema deberá mostrar eventos y alertas mediante un dashboard web.
\item \textbf{RF-07}: El sistema deberá mostrar el estado de los sensores monitorizados.
\item \textbf{RF-08}: El sistema deberá permitir la visualización de estadísticas de funcionamiento.
\item \textbf{RF-09}: El sistema deberá clasificar las anomalías según su nivel de riesgo.
\item \textbf{RF-10}: El sistema deberá permitir la configuración de determinados parámetros operativos.
\end{itemize}

\subsection{Requisitos no funcionales}

\begin{itemize}
\item \textbf{RNF-01}: El sistema deberá ser modular y fácilmente mantenible.
\item \textbf{RNF-02}: El sistema deberá utilizar tecnologías de código abierto.
\item \textbf{RNF-03}: El sistema deberá ejecutarse con un consumo reducido de recursos.
\item \textbf{RNF-04}: El sistema deberá ofrecer una interfaz sencilla e intuitiva.
\item \textbf{RNF-05}: El sistema deberá facilitar futuras ampliaciones funcionales.
\end{itemize}

\section{Especificación de requisitos}
CU-01 Consultar dashboard
CU-02 Consultar eventos
CU-03 Consultar alertas
CU-04 Consultar anomalías
CU-05 Consultar estado de sensores
CU-06 Configurar parámetros del sistema

